\documentclass[a4paper]{article}
\usepackage[margin=2.5cm]{geometry} % article's default margins: ~4 cm
\usepackage{titling}
\setlength{\droptitle}{-2cm} % \maketitle leaves a lot of space above the title
\usepackage{amsmath}
\usepackage[md]{juliatex} % [md]: also analysis.md (pandoc)
\makeatletter
\setlength{\@fptop}{0pt} % a page with only figures starts at the top, not centered
\makeatother

\title{Company Valuation with CorporateFinance.jl}
\author{Miguel Ferrera}

\begin{document}
\maketitle

\section{Cost of capital}

The target company has no traded shares, so its beta comes from a comparable company:
we unlever the comparable's beta and relever it to the target's capital structure.
The cost of equity follows from CAPM with a country risk premium, and the discount rate is the WACC.

\begin{juliacode}
using CorporateFinance
bu = unlever_beta(beta = 1.3, tax = 0.30, E = 500_000, D = 500_000)
bl = relever_beta(beta = bu, tax = 0.30, E = 800_000, D = 200_000)
ke = capm(rf = 0.04, beta = bl, erp = 0.055, crp = 0.02)
w = wacc(ke = ke, kd = 0.06, tax = 0.30, E = 800_000, D = 200_000)
map(v -> round(v, digits = 4), (beta_unlevered = bu, beta = bl, ke = ke, wacc = w))
\end{juliacode}

The cost of equity is \juliav{round(100ke, digits = 2)}\,\% and the WACC is \juliav{round(100w, digits = 2)}\,\%.

\section{Discounted cash flows}

Five years of free cash flow are projected, followed by a Gordon terminal value with 2\,\% long-term growth.
Net debt (debt minus cash) takes us from the enterprise value to the equity value.

\begin{juliacode}
fcf = [120_000, 135_000, 150_000, 160_000, 170_000]
r = dcf_valuation(fcf; wacc = w, g = 0.02)
equity = equity_value(ev = r.enterprise_value, debt = 200_000, cash = 30_000)
map(v -> round(Int, v), (pv_fcf = r.pv_fcf, pv_terminal = r.pv_terminal,
                         enterprise_value = r.enterprise_value, equity = equity))
\end{juliacode}

\begin{juliarun}
using Plots
thousands(v) = replace(string(round(Int, v)), r"(?<=\d)(?=(\d{3})+$)" => ",")

# Valuation bridge (thousands): discounted flows + terminal value = EV; EV - net debt = equity
steps = [fcf ./ (1 + w) .^ (1:5); r.pv_terminal] ./ 1000
top = cumsum(steps)
ev, eq = r.enterprise_value / 1000, equity / 1000
lo = [top .- steps; 0; eq; 0]
hi = [top; ev; ev; eq]
names = ["Year 1", "Year 2", "Year 3", "Year 4", "Year 5", "Terminal", "EV", "Net debt", "Equity"]
labels = [thousands.(steps); thousands(ev); "-" * thousands(ev - eq); thousands(eq)]
colors = [fill(:steelblue, 5); :darkorange; :gray50; :firebrick; :seagreen]

function bridge(k) # shows the first k bars
    title = k <= 7 ? "Enterprise value: $(thousands(1000(k == 0 ? 0 : hi[min(k, 7)])))" :
                     "Equity value: $(thousands(1000eq))"
    p = plot(legend = false, size = (720, 400), title = title, ylabel = "Thousands",
             ylims = (0, 1.15ev), xlims = (0.4, 9.6), xticks = (1:9, names),
             xtickfontsize = 8, grid = :y, yformatter = thousands)
    for i in 1:k
        plot!(p, Shape([i - 0.36, i + 0.36, i + 0.36, i - 0.36], [lo[i], lo[i], hi[i], hi[i]]),
              color = colors[i], linecolor = colors[i])
        annotate!(p, i, hi[i] + 0.03ev, text(labels[i], 8))
    end
    p
end

mkpath("resources")
savefig(bridge(9), "resources/bridge.pdf")
savefig(bridge(9), "resources/bridge.svg") # for analysis.md (GitHub does not show PDF)
# Animation for the README (a PDF cannot show GIF): the bridge is built step by step
anim = @animate for k in [0:9; fill(9, 6)]
    bridge(k)
end
gif(anim, "resources/bridge.gif", fps = 2)
\end{juliarun}

\begin{figure}[htbp]
  \centering
  \juliagraphics[width=\linewidth]{resources/bridge.pdf}
  \caption{Valuation bridge, in thousands, generated by Julia at compile time.}%
  \label{fig:bridge}
\end{figure}

As figure~\ref{fig:bridge} shows, the enterprise value is \juliav{thousands(r.enterprise_value)} and the equity value is \juliav{thousands(equity)}.
The terminal value accounts for \juliav{round(Int, 100r.pv_terminal / r.enterprise_value)}\,\% of the enterprise value, which is typical of a five-year forecast.

\section{Sensitivity}

Because the terminal value weighs so much, the result depends on the discount rate and on the long-term growth.
The table shows the enterprise value, in thousands, for WACC one point below and above the base case (rows) and for growth between 1.5\,\% and 2.5\,\% (columns).

\begin{juliacode}
waccs = round.(w .+ (-0.01:0.01:0.01), digits = 4)
gs = [0.015, 0.02, 0.025]
round.(Int, dcf_sensitivity(fcf; waccs, gs) ./ 1000)
\end{juliacode}

\end{document}
